\section{Feature comparison with OpenRocket}\label{app:or-features}

\Cref{app:or-context} sets the terms of this comparison: OpenRocket is the baseline because it
is the de facto standard for what a mature open-source model-rocket simulator offers, and the
comparison measures distance, not conformance. This appendix is the inventory behind that
judgment --- a matrix over sixteen feature areas: the components catalog; mass and inertia
modeling; CP and stability aerodynamics; drag decomposition; degrees of freedom and flight
dynamics; atmosphere, gravity, and wind; recovery; staging, clusters, and flight configurations;
motor data and selection; the design UI; 3D visualization; plotting and analysis; file formats
and interoperability; extensibility and scripting; optimization and dispersion; and
documentation, localization, platform, and distribution. The areas and the QtRocket verdicts
follow the companion architecture review's matrix; the OpenRocket column has been refreshed
against the source research gathered for \cref{app:or-methods,app:or-architecture}, and where
the older review and the newer research disagreed, the research prevails and the claim is stated
conservatively.

Two evidentiary standards apply, one per column, and the asymmetry is deliberate. QtRocket
claims carry \code{file:line} citations into the pinned commit \code{254cd41} and are grounded
in the body sections
(\cref{sec:parts,sec:motors,sec:simcore,sec:environment,sec:aero,sec:interfaces,sec:persistence,sec:testing}).
OpenRocket claims are condensed from the project's documentation, release notes, and source tree
at tag \code{release-24.12}, cited by footnote URL; they were checked against the repository in
July 2026 but are not held to the file-and-line standard of the QtRocket
column.\footnote{Primary sources: \url{https://openrocket.info};
\url{https://openrocket.readthedocs.io/en/latest/}; the \code{openrocket/openrocket} repository
at tag \code{release-24.12} (\url{https://github.com/openrocket/openrocket}); and the technical
documentation at \url{https://openrocket.sourceforge.net/techdoc.pdf}.} Version pinning matters
on both sides: 24.12 is the final stable GitHub release as of this writing --- the beta shipped
December 2024 and the final release is dated July~27, 2025\footnote{GitHub releases of
\code{openrocket/openrocket}: \code{release-24.12} published 2025-07-27;
\url{https://github.com/openrocket/openrocket/releases}.} --- while development continues on the
\code{unstable} branch. The matrix names \code{unstable} explicitly wherever a claim would
otherwise mislead.

The status markers grade the QtRocket column only; the OpenRocket column describes.
\fpresent{} means implemented and consumed by production code. \fpartial{} means built but
incomplete or unconsumed --- and in this codebase that marker deserves a second look, because
several \fpartial{} verdicts cover machinery that is finished and unit-tested but wired to
nothing (\cref{sec:incomplete}); the distance those rows measure is a consumption gap, not a
construction gap. \fabsent{} means no trace at the pinned commit, each such verdict grounded in
a failed search of the tree rather than an impression. \fwip{} means under development on an
unmerged branch, reported as such and never silently folded into the pinned-commit picture.

\subsection{The sixteen-area matrix}\label{app:or-features-matrix}

\begingroup
\footnotesize
\setlength{\tabcolsep}{5pt}
\begin{longtable}{@{}>{\raggedright\arraybackslash}p{0.92in}>{\raggedright\arraybackslash}p{2.32in}>{\raggedright\arraybackslash}p{2.84in}@{}}
\caption{Feature-area comparison: OpenRocket 24.12 versus QtRocket at commit
\code{254cd41}.}\label{tab:or-features-matrix}\\
\toprule
\textbf{Area} & \textbf{OpenRocket} & \textbf{QtRocket} \\
\midrule
\endfirsthead
\multicolumn{3}{@{}l}{\textit{\tablename~\thetable{} (continued)}}\\[2pt]
\toprule
\textbf{Area} & \textbf{OpenRocket} & \textbf{QtRocket} \\
\midrule
\endhead
\midrule
\multicolumn{3}{r@{}}{\textit{continued on next page}}\\
\endfoot
\bottomrule
\endlastfoot
Rocket components (part catalog) &
Roughly two dozen concrete types under one abstract \code{RocketComponent} tree: six nose-cone shape
families, transitions, inner tubes, couplers, centering rings, bulkheads, lugs, rail buttons,
four fin classes including free-form, recovery devices, mass objects, pods, boosters. Per-class
\code{isCompatible} child rules; per-component mass/CG/$C_d$ overrides; a commercial preset
database. &
\fpartial{} Five types (\src{model/parts/PartFactory.cpp}{32}); geometry immutable after
construction, so live dimension editing has no mutation path. Distinctive: intent-based
\code{StationLink} placement (\src{model/parts/Placement.h}{49}) whose overlap sweep hard-fails
self-intersecting designs (\src{model/parts/Part.cpp}{184}) --- stronger than warn-and-simulate
(\cref{sec:placement}). \\
\addlinespace
Mass / inertia modeling &
Mass from volume and density per component; an external \code{MassCalculator} walks the tree
each force evaluation; two scalar moments (longitudinal, axial) --- the tensor is treated as
diagonal and axisymmetric. Overrides: overridden mass rescales inertia linearly, overridden CG
shifts it by parallel axis. Motor mass and CG interpolated between curve samples during burn. &
\fpresent{} Lazy two-pass time-aware composition with parallel-axis shifts
(\src{model/parts/Part.cpp}{175}); mass-delta cache freezes bit-stably after burnout
(\src{model/parts/Part.cpp}{130}); closed-form full $3\times3$ tensors pinned by numeric oracles
(\src{model/tests/InertiaTensorsTests.cpp}{186}). The $R\,I\,R^{T}$ child-tensor rotation is
omitted pending 6-DOF (\src{model/parts/Part.cpp}{217}). \\
\addlinespace
Aerodynamics: CP / stability &
Extended Barrowman recomputed every force evaluation at the instantaneous angle of attack and
Mach number ($C_{N_\alpha} \equiv C_N/\alpha$): Galejs body lift, free-form fin MAC integrals,
fin--body interference, pitch damping, roll forcing and damping; live CP/CG markers refreshed by
a background simulation while editing. &
\fpartial{} Per-part Barrowman math exists and is unit-tested, with composite CP on a shared tip
datum (\src{sim/Aero.h}{25}, \src{model/parts/Part.cpp}{231}) --- but zero production consumers,
and no surface shows CP, CG, or margin. Latent defect: mirrored \code{FinSet} chordwise sign
convention (\cref{sec:aero}). \\
\addlinespace
Aerodynamics: drag decomposition &
Component build-up: roughness-aware fully-turbulent skin friction with a roughness-limited
floor; per-shape transonic/supersonic nose pressure drag from experimental curves; fin
leading/trailing-edge terms; base drag with motor area subtracted under thrust; parasitic lug
drag; axial drag scaled empirically with angle of attack. &
\fabsent{} One user-set scalar $C_d$ (default 1.0) times one reference area
(\src{model/RocketModel.cpp}{81}); every part's drag contribution returns zero
(\src{model/parts/BodyTube.cpp}{41}); no Mach or Reynolds number exists.
\code{deriveReferenceAreaFromGeometry()} is correct but uncalled
(\src{model/RocketModel.cpp}{48}), so all 25 design fixtures under \code{tests/data/designs/}
fly on a stale or default 38\,mm disc. \\
\addlinespace
Degrees of freedom / flight dynamics &
Quaternion 6-DOF ascent over a two-scalar inertia model, degrading deliberately to 3-DOF for
recovery descent and tumbling stages; RK4 with a-priori heuristic step limits, no embedded error
pair; validated against flight and wind-tunnel data (headline accuracy roughly 10\% at subsonic
speeds). &
\fpartial{} 3-DOF point mass, heavily tested end-to-end. Thrust hardwired along world $+z$
(\src{model/RocketModel.cpp}{70}); \code{getTorques()} returns zero with no caller
(\src{model/RocketModel.cpp}{94}); orientation integrator commented out
(\src{sim/Propagator.h}{103}). Ahead: embedded-error adaptive RKF45
(\src{sim/RK45Solver.h}{75}); typed termination taxonomy (\src{sim/Propagator.h}{35}). \\
\addlinespace
Atmosphere / gravity / wind &
Eight-layer ISA with launch-site overrides; pink-noise turbulence (two-pole Kasdin filter,
$\alpha = 5/3$) and, new in 24.12, multi-level winds with CSV import; flat, spherical, or
WGS84--Vincenty geodesy with Coriolis; Somigliana latitude-dependent gravity; launch rod/rail. &
\fpartial{} Three atmospheres (US Standard 1976 validated to 0.1--0.5\% up to 80\,km) and true
inverse-square spherical gravity (\src{sim/SphericalGravityModel.cpp}{28}). Wind is dead code:
returns $(0,0,0)$, zero call sites (\src{sim/WindModel.cpp}{16}); no rod/rail (pad-hold clamp
only, \src{sim/Propagator.cpp}{105}), no launch site, no Coriolis. \\
\addlinespace
Recovery &
Event-driven deployment over sixteen \code{FlightEvent} types\footnote{\code{FlightEvent.java}
under \code{core/src/main/java/info/openrocket/core/simulation/}: \code{LAUNCH},
\code{IGNITION}, \code{LIFTOFF}, \code{LAUNCHROD}, \code{BURNOUT}, \code{EJECTION\_CHARGE},
\code{STAGE\_SEPARATION}, \code{APOGEE}, \code{RECOVERY\_DEVICE\_DEPLOYMENT}, \code{GROUND\_HIT},
\code{SIMULATION\_END}, \code{ALTITUDE}, \code{TUMBLE}, \code{SIM\_WARN}, \code{SIM\_ABORT},
\code{EXCEPTION}.} with per-trigger delays; dual deployment via altitude events; parachute $C_d$
default 0.8 and an empirical streamer drag law; dedicated 3-DOF descent and tumble steppers with
exact apogee/ground root solves. &
\fabsent{} Nothing changes at apogee; flight ends ballistically at the $z<0$-descending
terminate (\src{model/RocketModel.cpp}{62}). Ejection delays are parsed and persisted but
consumed by no flight logic; apogee exists only as a post-hoc statistic. \\
\addlinespace
Staging / clusters / flight configurations &
Unlimited serial stages plus \code{ParallelStage} boosters and pods; stage separation forks the
simulation, one \code{FlightDataBranch} per branch; motor clusters via \code{Clusterable} inner
tubes; named flight configurations keyed by UUID with per-stage active flags. &
\fabsent{} Single motor, single stage: first-match DFS (\src{model/RocketModel.cpp}{13});
\code{Stage.h} is a commented-out ``Not yet'' include (\src{model/RocketModel.h}{22}). Ignition
timing is half-plumbed (\code{ThrustCurve::setIgnitionTime} has zero call sites); the motor sits
at the airframe CM with a hardcoded zero offset (\src{model/RocketModel.cpp}{124}). \\
\addlinespace
Motor data \& selection &
Complete thrustcurve.org corpus bundled offline: 24.12 ships a serialized snapshot; the
\code{unstable} branch replaces it with an SQLite database published from a separate
motor-database repository (1458 motors, 1591 curves at the checked
snapshot).\footnote{\url{https://github.com/openrocket/motor-database}; on \code{unstable} the
bundled \code{datafiles/thrustcurves/initial\_motors.db} carries metadata
\code{motor\_count:~1458}, \code{curve\_count:~1591} and a published update URL.} RASP
\code{.eng} and RockSim \code{.rse} user files; rich chooser with search, filters, curve plots,
delay presets. &
\fpartial{} Three ingest paths behind one facade: RSE (\src{model/RSEDatabaseLoader.cpp}{20}),
strict RASP (\src{model/RASPLoader.cpp}{253}), live thrustcurve.org search behind a
fake-injectable seam (\src{model/ThrustCurveClient.h}{68}) --- live in-app search the baseline
lacks. Ships with an empty database (the one bundled \code{.rse} is loaded only by tests);
per-sample mass/CG data discarded (\src{model/RSEDatabaseLoader.cpp}{89}); common-name keying
replaces on collision --- last write wins, distinguished only in the logs. \\
\addlinespace
Design UI &
Component tree with per-component config dialogs, drag-and-drop, and undo; live 2D schematic
with CG/CP markers; a background-simulation thread pool re-runs the flight on every edit to keep
the altitude and stability readouts live.\footnote{\code{RocketPanel.java}
(\code{swing/src/main/java/info/openrocket/swing/gui/scalefigure/}) owns a daemon
background-simulation executor; every design edit triggers \code{updateExtras()}, which re-runs
a simulation to refresh the on-canvas readouts.} &
\fwip{} At the pinned commit, a motor-and-launch front end only: design widgets are constructor
stubs (\src{gui/RocketTreeView.cpp}{3}), File actions ship disabled
(\src{gui/MainWindow.ui}{117}), authoring is CLI-only (\src{cli/Repl.cpp}{784}). The unmerged
\code{RocketTreeViewImplementation} branch adds a read-only tree and \code{.qrd} open/save; no
add/edit UI, no property panel, no schematic (\cref{sec:interfaces}). \\
\addlinespace
3D visualization &
Embedded JOGL 3D finished/unfinished views; per-component appearance, textures, and decals;
PhotoStudio photorealistic rendering. &
\fpresent{} Standalone \code{qtrocket-visualizer} meshes the tree
(\src{visualizer/RocketMesh.cpp}{404}) and positions meshes with the simulator's own placement
resolver (\src{visualizer/RocketMesh.cpp}{433}), so rendered and simulated geometry cannot
diverge; overlap offenders render in error red (\src{visualizer/RocketGLWidget.cpp}{617}). Not
embedded in the GUI; motors render no geometry; no tests. \\
\addlinespace
Plotting / analysis &
Fifty-plus typed flight-data series on configurable axes with event overlays; user-defined
custom expressions; component analysis; CSV export with selectable columns. &
\fpartial{} Three fixed plots (\src{gui/AnalysisWindow.cpp}{36}); CLI launch summary plus a
fixed-schema CSV auto-written to a hardcoded \code{./qtrocket\_run.csv}
(\src{cli/Repl.cpp}{104}). No configurable plots, no event annotations, no stability outputs. \\
\addlinespace
File formats \& interop &
\code{.ork} ZIP-plus-XML with a content-sniffing loader (ZIP, gzip, and bare XML all load) and
format versions back to 1.0; RockSim \code{.rkt} import and export; RASAero \code{.CDX1} import;
OBJ, SVG, and CSV export. &
\fpartial{} Versioned \code{.qrd} XML round-trips the tree with placement intent, re-resolved on
load (\src{model/DesignSerializer.cpp}{93}); fail-safe atomic loading
(\src{model/DesignSerializer.cpp}{255}). Legacy 0.1 files rejected with no migration path
(\src{model/DesignSerializer.cpp}{194}); mass overrides lost on save; no foreign import or
export of any kind. \\
\addlinespace
Simulation extensibility / scripting &
Three listener interfaces --- lifecycle, event veto, and paired pre/post hooks around every
physics sub-model (gravity, wind, thrust, aero, mass) that can override results per step;
shipped simulation extensions; JavaScript scripting via GraalJS; annotation-scanned plugin JARs. &
\fpartial{} No listener, plugin, or scripting mechanism. Partial substitute the baseline lacks:
a 33-command scriptable CLI over a machine-parseable OK/ERR protocol
(\src{cli/Repl.cpp}{232}), weakened as automation by an always-zero exit code. \\
\addlinespace
Optimization / dispersion &
Rocket Optimizer: parameter modifiers, altitude/velocity/stability goals, and stability-domain
constraints evaluated over parallel simulation runs. Monte Carlo dispersion absent, on the
project's own roadmap. &
\fabsent{} Neither optimizer nor dispersion exists (greps return zero hits); since the baseline
also lacks Monte Carlo, the gap is the optimizer only. The deterministic core and scriptable CLI
make external sweeps feasible today, but nothing in-product supports them. \\
\addlinespace
Docs / i18n / platform / distribution &
GPLv3 Java~17 split into JPMS core and Swing modules; installers with bundled JRE for
Windows/macOS/Linux plus a Snap package; localization into roughly fifteen languages; published
technical documentation with a validation chapter. &
\fpartial{} Strong engineering infrastructure: five-platform CI with \code{-Werror}
(\src{.github/workflows/cmake-multi-platform.yml}{28}), \code{llvm-cov} coverage, 239 tests.
Distribution essentially absent: one install rule (\src{CMakeLists.txt}{308}), no CPack or
installers; i18n is dead scaffolding (\src{CMakeLists.txt}{201}), English only. \\
\end{longtable}
\endgroup

\subsection{Where QtRocket is ahead}\label{app:or-features-ahead}

Fairness cuts both ways. Sixteen rows of mostly \fpartial{} and \fabsent{} would misrepresent
the comparison if the appendix did not also record the places where QtRocket offers something
the baseline does not. Seven are defensible at the pinned commit.

\begin{description}
\item[Live motor search.] Both front ends can search thrustcurve.org interactively --- facet
metadata, filtered search, per-result curve download --- through one client behind a
fake-injectable seam (\src{model/ThrustCurveClient.h}{68}). OpenRocket's chooser is
offline-first, working against a corpus bundled at release time (the 24.12 serialized snapshot;
on \code{unstable}, the SQLite database above). The trade cuts both ways --- OpenRocket is
complete out of the box where QtRocket starts empty (\cref{sec:motors}) --- but interactive
search against the canonical registry is a capability the baseline's user surface does not
offer.

\item[Embedded-error adaptive integration.] QtRocket's RKF45 is a true embedded pair: fourth-
and fifth-order estimates from shared stages, a per-step error norm, step rejection and
re-integration, and a growth clamp pinned by a regression test (\src{sim/RK45Solver.h}{75}).
OpenRocket integrates ascent with RK4 whose step is chosen a priori as the minimum of heuristic
caps --- pitch-angle and roll-rate limits, event distance, a $1.5\times$ growth limiter ---
computed after evaluating only the first stage; the methods research confirmed there is no error
estimate and no step rejection anywhere in the stepper\footnote{\code{RK4SimulationStepper.java}
under \code{core/src/main/java/info/openrocket/core/simulation/}; the stepper ships at tag
\code{release-24.12}, and the eight step-size candidates --- none involving an error
estimate --- were verified in the class's code comments on the \code{unstable} development
branch (fetched 2026-07). See \cref{app:or-methods}.} (\cref{app:or-methods}). In fairness, those heuristics are
attitude-aware and tuned to 6-DOF oscillation, a problem QtRocket does not yet face.

\item[A hang-proof, typed termination taxonomy.] Every exit from the run loop is one of five
typed verdicts --- nominal, no-liftoff, non-finite state, time/iteration cap, integrator error
(\src{sim/Propagator.h}{35}) --- and the loop is bounded by both a simulated-time cap and an
iteration cap, so it cannot hang. The CLI consumes the enum in an exhaustive switch with no
default, so a new verdict fails to compile until every handler is written
(\src{cli/Repl.cpp}{41}). OpenRocket reports aborts and warnings as flight events
(\code{SIM\_ABORT}, \code{SIM\_WARN}) in the flight data; QtRocket's distinction is the
exhaustive partition and the compile-time obligation it imposes on consumers.

\item[A scriptable headless REPL.] Thirty-three commands, exactly one machine-parseable
\code{OK}/\code{ERR} line per command, and designs authored, persisted, and flown entirely from
a pipe (\src{cli/Repl.cpp}{232}). OpenRocket's core module is deliberately headless and
reusable, but the project ships no interactive console: scripted use means writing Java against
the module or using the separately maintained Python
bindings.\footnote{\url{https://github.com/openrocket/orhelper}.}

\item[One resolver for physics and rendering.] The visualizer positions every mesh from the same
\code{resolvePlacements} output the mass and inertia composition consumes
(\src{visualizer/RocketMesh.cpp}{433}, \src{model/parts/Placement.h}{148}), so rendered and
simulated geometry cannot diverge by construction. OpenRocket's 2D figure and 3D views are
mature, separate rendering paths over the component tree; nothing makes their geometry and the
simulator's identical by construction.

\item[Fail-closed geometry diagnostics.] A self-intersecting design is a hard, located failure
before any number is produced: the composite mass computation refuses with the first overlap
diagnostic (\src{model/parts/Part.cpp}{184}), and the visualizer paints the offender in error
red from the same cached verdict (\src{visualizer/RocketGLWidget.cpp}{617}). OpenRocket warns
and simulates anyway --- more forgiving, but a produced number does not guarantee a coherent
model.

\item[True inverse-square gravity.] The spherical gravity option evaluates the central field
$-GM\,\mathbf{r}/|\mathbf{r}|^{3}$ at the integrator's trial position
(\src{sim/SphericalGravityModel.cpp}{28}). OpenRocket's \code{WGSGravityModel} applies
Somigliana latitude-dependent surface gravity with a spherical-approximation
$(R/(R+h))^{2}$ altitude correction.\footnote{\code{WGSGravityModel.java} under
\code{core/src/main/java/info/openrocket/core/models/gravity/}.} At model-rocket altitudes the
numbers are close, and OpenRocket's latitude dependence is the more useful refinement there; the
QtRocket option's distinction is being an exact central vector field rather than a corrected
scalar magnitude --- the right shape for the high-altitude regime its US Standard Atmosphere
already covers.
\end{description}

None of this reverses the overall verdict. OpenRocket is a mature product with a broad surface
and a published validation record; QtRocket is a much smaller simulator whose engineering
discipline runs ahead of its feature count. The matrix should be read with
\cref{sec:incomplete} in hand: much of its \fpartial{} column is finished machinery awaiting a
consumer, so the distance it measures is overwhelmingly a consumption gap, not a construction
gap.
